\chapter{排出废物}\label{ch:g7:removing-wastes}

这一年的主题大多是“输送”：\omterm{def:g7:what-is-respiration:respiration}{氧气}送入，\omterm{def:g7:digestion-and-nutrients:families}{营养物质}送入，
燃料送到每个工作的\omterm{def:g5:first-look-at-cells:cell}{细胞}。但任何车间都会产生废物，
一家从不收垃圾的运送公司，
很快就会污染自己所在的城市。\omterm{def:g7:what-is-respiration:respiration}{二氧化碳}的去路已定——
在\omterm{def:g7:lungs-and-gas-exchange:alveoli}{肺泡}处呼出。
本章是全书的收尾，追踪的是\emph{其他}废物，
它们由身体里最低调的器官——\omterm{def:g7:removing-wastes:kidneys}{肾脏}——处理。

\section{二氧化碳之外的废物}

\begin{proposition}[细胞的其他“垃圾”]\label{prop:g7:removing-wastes:rubbish}
\omterm{def:g5:first-look-at-cells:cell}{细胞}的燃烧与构建
（\cref{prop:g7:what-is-respiration:power}，
\cref{prop:g6:matter-from-food:transform}）除了产生\omterm{def:g7:what-is-respiration:respiration}{二氧化碳}，
还会产生其他废物。尤其是利用构建物质的\omterm{def:g7:digestion-and-nutrients:families}{营养物质}时，
会留下一种含氮废物——
\emph{尿素}\index{尿素}——\omterm{prop:g4:journey-of-food:absorb}{血液}
从每个组织中收集它，方式与收集\omterm{def:g7:what-is-respiration:respiration}{二氧化碳}完全相同。\omterm{prop:g7:removing-wastes:rubbish}{尿素}
经过时无害，堆积起来则有毒，因此
必须持续排出体外。
\end{proposition}
\admitted

\begin{example}[两次收集，同一股血液]\label{ex:g7:removing-wastes:rounds}
\cref{ex:g7:heart-and-circulation:tour}中的同一滴血
同时进行两次收集：在大腿的\omterm{def:g7:heart-and-circulation:vessels}{毛细血管}处，
它既带上\omterm{def:g7:what-is-respiration:respiration}{二氧化碳}，\emph{也}带上\omterm{prop:g7:removing-wastes:rubbish}{尿素}。第一种在肺部卸下，
进入肺循环；第二种继续随\omterm{prop:g4:journey-of-food:absorb}{血液}流动——
流经肺部时，肺无法处理它——
直到体循环把它送往专为这类货物而设的器官，
即\omterm{def:g7:removing-wastes:kidneys}{肾脏}。
\end{example}

\section{肾脏，血液的过滤器}

\begin{definition}[肾脏与泌尿系统]\label{def:g7:removing-wastes:kidneys}
两个\emph{肾脏}\index{肾脏}——拳头大小，
形如豆子，位于背部偏低处——是清洁\omterm{prop:g4:journey-of-food:absorb}{血液}的
过滤器。每个肾脏都直接从体循环的主干道接收一条粗大的\omterm{def:g7:heart-and-circulation:vessels}{动脉}。
\omterm{prop:g4:journey-of-food:absorb}{血液}在肾脏内部，
经约一百万个微观的过滤单位过滤；
清洁后的\omterm{prop:g4:journey-of-food:absorb}{血液}由肾\omterm{def:g7:heart-and-circulation:vessels}{静脉}流回，
被移除的负担——\omterm{prop:g7:removing-wastes:rubbish}{尿素}、多余的水和多余的盐——
以\emph{尿液}\index{尿液}的形式排出，
沿两条输尿管流入\emph{膀胱}\index{膀胱}，
膀胱是可随时排空的储存器官。
\end{definition}

\begin{proposition}[先过滤，后回收]\label{prop:g7:removing-wastes:filter}
\omterm{def:g7:removing-wastes:kidneys}{肾脏}的诀窍是先大量滤出，再节约回收。
过滤单位首先将\omterm{prop:g4:journey-of-food:absorb}{血液}中的大量水分部分，
每天可达数十升，通过其精细的筛网，
\omterm{prop:g7:removing-wastes:rubbish}{尿素}、糖类、盐类一并滤出。随后，
沿每个单位的弯曲管道，身体\emph{回收}它需要的物质：
几乎全部的水、全部的\omterm{prop:g7:organs-at-work:bill}{葡萄糖}，
以及适量的盐，而\omterm{prop:g7:removing-wastes:rubbish}{尿素}和多余的物质
则留在滤液中。在滤出的数十升液体中，约有一升半
以\omterm{def:g7:removing-wastes:kidneys}{尿液}的形式排出。全部滤出，
取回有用的部分，废物便无处可藏。
\end{proposition}
\admitted

\begin{omfigure}
\begin{tikzpicture}[scale=0.9]
  % blood in
  \draw[->, very thick, omThm] (-0.2,3.0) -- (2.2,3.0);
  \node[font=\footnotesize, omThm, above, align=center] at (1.0,3.1)
    {血液流入 \\（含尿素）};
  % kidney box
  \draw[thick, fill=omDef!6, rounded corners=8pt]
    (2.3,1.0) rectangle (8.1,4.6);
  \node[font=\small\bfseries, omDef] at (5.2,4.25) {肾脏};
  % filter stage
  \draw[thick, omDef, fill=white] (2.7,2.2) rectangle (4.5,3.6);
  \node[font=\footnotesize, align=center] at (3.6,2.9)
    {过滤：\\ 全部滤出};
  % reclaim stage
  \draw[thick, omMeth, fill=white] (5.0,2.2) rectangle (7.7,3.6);
  \node[font=\footnotesize, align=center] at (6.35,2.9)
    {回收：\\ 水、葡萄糖、\\ 适量的盐};
  \draw[->, thick] (4.5,2.9) -- (5.0,2.9);
  % blood out
  \draw[->, very thick, omDef] (8.1,3.0) -- (10.5,3.0);
  \node[font=\footnotesize, omDef, above, align=center] at (9.3,3.1)
    {清洁后的 \\ 血液流出};
  % urine out
  \draw[->, very thick, omProp] (6.35,2.2) -- (6.35,0.0);
  \node[font=\footnotesize, omProp, right, align=left] at (6.5,0.55)
    {尿液：尿素、多余的 \\ 水和盐};
  \node[font=\footnotesize, align=center] at (6.35,-0.35)
    {流向膀胱};
\end{tikzpicture}

{\small \omterm{def:g7:removing-wastes:kidneys}{肾脏}的双阶段诀窍：大量滤出，节约回收。
未被回收的部分——\omterm{prop:g7:removing-wastes:rubbish}{尿素}和多余的物质——
即为\omterm{def:g7:removing-wastes:kidneys}{尿液}。}
\end{omfigure}

\begin{example}[从节约中读出道理]\label{ex:g7:removing-wastes:thrift}
回收阶段可以解释日常中的观察结果。
喝水多时，\omterm{def:g7:removing-wastes:kidneys}{肾脏}回收的水分较少：尿量多，
颜色浅。在炎热的比赛中大量出汗，
\omterm{def:g7:removing-wastes:kidneys}{肾脏}几乎全部回收：尿量少，颜色深——身体
在保护自己的水分
（\cref{prop:g4:what-plants-need:needs}讲过水对\omterm{def:g1:plants-around-us:plant}{植物}的价值；
\omterm{def:g7:removing-wastes:kidneys}{肾脏}则是\omterm{def:g1:animals-around-us:animal}{动物}一侧的同一套
经济原则）。\omterm{def:g7:removing-wastes:kidneys}{尿液}中出现\omterm{prop:g7:organs-at-work:bill}{葡萄糖}，过滤环节放行了它，
而回收本应将其全部取回；
这是一个典型的医学信号，说明某个环节——
回收或血糖——出了问题。
\end{example}

\begin{remark}[皮肤也来帮忙]\label{rem:g7:removing-wastes:sweat}
\emph{汗液}\index{汗液}是同一门生意的
副业：水和盐（含微量\omterm{prop:g7:removing-wastes:rubbish}{尿素}）经\omterm{def:g1:the-five-senses:sense}{皮肤}
排出——不过它的首要任务是
\cref{ex:g7:organs-at-work:felt}所说的散热。肺、
\omterm{def:g7:removing-wastes:kidneys}{肾脏}和\omterm{def:g1:the-five-senses:sense}{皮肤}分担了全身的废物排出：
\omterm{def:g7:what-is-respiration:respiration}{二氧化碳}交给前者，\omterm{prop:g7:removing-wastes:rubbish}{尿素}和多余物质交给
后者，热量和少量含盐的汗水交给\omterm{def:g1:the-five-senses:sense}{皮肤}。
不会有任何毒物堆积——在健康的
身体里，废物永远不会占上风。
\end{remark}

\section{全年的机器，完整版}

\begin{proposition}[完整的供给与排出循环]\label{prop:g7:removing-wastes:complete}
有了\omterm{def:g7:removing-wastes:kidneys}{肾脏}，本年度的讲述便合成
一台完整的机器。\omterm{prop:g4:journey-of-food:absorb}{血液}在双循环中被泵送，
（\cref{def:g7:heart-and-circulation:double}），
在\omterm{def:g7:lungs-and-gas-exchange:alveoli}{肺泡}处装载\omterm{def:g7:what-is-respiration:respiration}{氧气}，在\omterm{prop:g7:digestion-and-nutrients:absorption}{绒毛}处装载\omterm{def:g7:digestion-and-nutrients:families}{营养物质}；
将两者送达每个工作的\omterm{def:g5:first-look-at-cells:cell}{细胞}，供燃烧与
构建之用；并收集废物——\omterm{def:g7:what-is-respiration:respiration}{二氧化碳}交给
肺，\omterm{prop:g7:removing-wastes:rubbish}{尿素}和多余物质交给\omterm{def:g7:removing-wastes:kidneys}{肾脏}。四年级和七年级
讲过的每个器官，都是这单一循环中的一站；
每次交换都要穿过一层\omterm{def:g7:heart-and-circulation:vessels}{毛细血管}壁，而整个
循环由\cref{prop:g7:what-is-respiration:power}所述的
慢火驱动，逐\omterm{def:g5:first-look-at-cells:cell}{细胞}进行。
\end{proposition}
\admitted

\begin{method}[全身审计]\label{met:g7:removing-wastes:audit}
对任一物质，审核其路径：
\begin{enumerate}
  \item 进入：口腔与\omterm{prop:g7:digestion-and-nutrients:absorption}{绒毛}，或鼻子与\omterm{def:g7:lungs-and-gas-exchange:alveoli}{肺泡}；
  \item 运输：经哪段循环，途经哪些特殊
        站点（\omterm{prop:g7:digestion-and-nutrients:absorption}{肝脏}）；
  \item 利用：被燃烧作为燃料，被用于构建，或两者皆非；
  \item 排出：经\omterm{def:g7:lungs-and-gas-exchange:alveoli}{肺泡}（气体）、\omterm{def:g7:removing-wastes:kidneys}{肾脏}（\omterm{prop:g7:removing-wastes:rubbish}{尿素}、多余物质）、
        \omterm{def:g1:the-five-senses:sense}{皮肤}（\omterm{rem:g7:removing-wastes:sweat}{汗液}的份额）——或者，若从未被吸收，
        则走残余物的路径
        （\cref{prop:g4:journey-of-food:absorb}）。
\end{enumerate}
对水、\omterm{prop:g7:organs-at-work:bill}{葡萄糖}和一片咸味薯片运行此审计：机器为每一步都
设有一站，这正是全年发现的要点，
浓缩于这一次审计之中。
\end{method}

\section{练习}

\begin{exercise}[$\star$]\label{exo:g7:removing-wastes:1}
\omterm{prop:g7:removing-wastes:rubbish}{尿素}是什么——它源于何种
过程，又为何必须排出？
\end{exercise}

\begin{exercise}[$\star$]\label{exo:g7:removing-wastes:2}
哪些废物经肺排出，哪些经\omterm{def:g7:removing-wastes:kidneys}{肾脏}排出，
哪些经\omterm{def:g1:the-five-senses:sense}{皮肤}排出？
\end{exercise}

\begin{exercise}[$\star$]\label{exo:g7:removing-wastes:3}
描述\omterm{def:g7:removing-wastes:kidneys}{肾脏}：数目、大小、位置，以及
它们与\omterm{def:g7:removing-wastes:kidneys}{膀胱}的管道连接。
\end{exercise}

\begin{exercise}[$\star$]\label{exo:g7:removing-wastes:4}
叙述“先过滤、后回收”的机制，并说明
框定它的两个体积量。
\end{exercise}

\begin{exercise}[$\star$]\label{exo:g7:removing-wastes:5}
哪些物质被完全回收，应测量什么，
哪些完全不被回收？
\end{exercise}

\begin{exercise}[$\star$]\label{exo:g7:removing-wastes:6}
为何\omterm{def:g7:removing-wastes:kidneys}{尿液}有时量多、颜色浅，
另一些时候却量少、颜色深？
\end{exercise}

\begin{exercise}[$\star\star$]\label{exo:g7:removing-wastes:7}
为何肺不能排出\omterm{prop:g7:removing-wastes:rubbish}{尿素}，而\omterm{def:g7:removing-wastes:kidneys}{肾脏}也不能
足够快地排出\omterm{def:g7:what-is-respiration:respiration}{二氧化碳}？分工
体现了两种废物怎样的形态差异？
\end{exercise}

\begin{exercise}[$\star\star$]\label{exo:g7:removing-wastes:8}
解释为何\omterm{def:g7:removing-wastes:kidneys}{尿液}中的\omterm{prop:g7:organs-at-work:bill}{葡萄糖}会引起医生警觉，
请结合过滤与回收两个阶段说明。
\end{exercise}

\begin{exercise}[$\star\star$]\label{exo:g7:removing-wastes:9}
追踪一单位\omterm{prop:g7:removing-wastes:rubbish}{尿素}，从大腿\omterm{def:g3:bones-and-muscles:muscle}{肌肉}的\omterm{def:g7:heart-and-circulation:vessels}{毛细血管}
出发，逐站追踪至\omterm{def:g7:removing-wastes:kidneys}{膀胱}。
\end{exercise}

\begin{exercise}[$\star\star$]\label{exo:g7:removing-wastes:10}
运用\cref{met:g7:removing-wastes:audit}，对
运动后饮用的一杯水进行审核。
\end{exercise}

\begin{exercise}[$\star\star$]\label{exo:g7:removing-wastes:11}
过滤功能失效的肾病患者要使用一种机器，
它借助人工膜为其清洁\omterm{prop:g4:journey-of-food:absorb}{血液}。
该机器模仿了\cref{prop:g7:removing-wastes:filter}
中的哪一个阶段？为何治疗须
每隔几天重复一次？
\end{exercise}

\begin{exercise}[$\star\star\star$]\label{exo:g7:removing-wastes:12}
写出全年的总结段落，
即\cref{prop:g7:removing-wastes:complete}所承诺的：
一条\omterm{prop:g4:journey-of-food:absorb}{血液}、两条循环、四道边界（\omterm{def:g7:lungs-and-gas-exchange:alveoli}{肺泡}、\omterm{prop:g7:digestion-and-nutrients:absorption}{绒毛}、
\omterm{def:g3:bones-and-muscles:muscle}{肌肉}\omterm{def:g7:heart-and-circulation:vessels}{毛细血管}、\omterm{def:g7:removing-wastes:kidneys}{肾脏}过滤器）——以及
整台机器所服务的慢火。
\end{exercise}

\section{问题：一个水分子的一天}

\begin{problem}[{周末问题——一杯水，
贯穿整台机器的审计}]\label{pb:g7:removing-wastes:1}
跟踪运动日早餐时喝下的一杯水，
将其作为全年的大审计。

\medskip\noindent\textbf{第一部分——进入与运输。}
\begin{enumerate}
  \item 水无需拆解
        （见\cref{def:g7:digestion-and-nutrients:families}）。
        说出它进入体内的途径，以及
        接收它的\omterm{def:g4:heart-and-blood:vessels}{血管}网络。
  \item 它的\omterm{prop:g4:journey-of-food:absorb}{血液}首先到达的\omterm{def:g3:sorting-living-things:criterion}{标准}第一站是哪个器官，
        以及即使水本身无需分类，为什么这一路径
        对\omterm{def:g7:digestion-and-nutrients:families}{营养物质}仍然合理？
  \item 到上午中段，水已遍布\omterm{prop:g4:journey-of-food:absorb}{血液}所到之处，
        说出它此时交替流经的两条循环
        的名称。
  \item 不久，其中一部分进入正在工作的\omterm{def:g5:first-look-at-cells:cell}{细胞}内，列出
        水在\omterm{def:g5:first-look-at-cells:cell}{细胞}内承担的两项作用（回想一下，
        \omterm{def:g5:first-look-at-cells:cell}{细胞}内的化学反应——如汁液中的反应、燃烧——
        都在湿润\omterm{def:g6:exploring-our-environment:components}{环境}中进行）。
\end{enumerate}

\medskip\noindent\textbf{第二部分——下午
的比赛。}
\begin{enumerate}[resume]
  \item 比赛期间，部分水通过一道边界离开身体，
        这道边界并非排泄器官。
        说出这一边界、它随水带出的其他物质，
        以及它当天的真正作用。
  \item 比赛结束后，这位运动员口渴，
        当晚她的尿量少且颜色深。请
        结合重吸收阶段解释这两种现象。
  \item 与此同时，比赛中劳累的\omterm{def:g3:bones-and-muscles:muscle}{肌肉}产生了额外的
        \omterm{def:g7:what-is-respiration:respiration}{二氧化碳}和\omterm{prop:g7:removing-wastes:rubbish}{尿素}。请为每种物质指定其离开
        身体的边界，以及把它送到那里的
        循环段。
  \item 尿量减少时，\omterm{def:g7:removing-wastes:kidneys}{肾脏}为何在整场比赛中仍
        持续工作？它们此时仍在清除
        哪些物质？
\end{enumerate}

\medskip\noindent\textbf{第三部分——审计结束。}
\begin{enumerate}[resume]
  \item 次日清晨，这杯水中有一部分已经随
        \omterm{def:g4:breathing:movements}{呼气}离开身体。请解释水如何经\omterm{def:g7:lungs-and-gas-exchange:alveoli}{肺泡}离开
        （见\cref{ex:g4:breathing:mirror}，
        几年前已说明过此点）。
  \item 请列出这杯水的完整账目：进入、
        储备与利用，以及它的四条排出途径——\omterm{def:g7:removing-wastes:kidneys}{尿液}、\omterm{rem:g7:removing-wastes:sweat}{汗液}、
        呼出的气，以及残渣中的少量水分。
  \item 审计表明，同一条\omterm{prop:g4:journey-of-food:absorb}{血液}依次为肺、\omterm{prop:g7:digestion-and-nutrients:absorption}{绒毛}、
        \omterm{def:g3:bones-and-muscles:muscle}{肌肉}、\omterm{def:g7:removing-wastes:kidneys}{肾脏}和\omterm{def:g1:the-five-senses:sense}{皮肤}服务。请用两句话说明，
        为何这种单一网络的设计，使每个器官的健康
        与其他所有器官的
        状况息息相关。
  \item 为全年收尾：哪一个单一概念——请说出它——
        本次审计的每一站都服务于它？
        （见\cref{prop:g7:what-is-respiration:power}，
        其中给出了答案。）
\end{enumerate}
\end{problem}
